为了直观展示本工作提出的方法与经典基于多边形、一阶 SDF 模型的架构差异，图 \ref{fig:flowchart} 对比了三种不同的接触动力学计算管线。

\begin{figure}[htbp]
    \centering
    \includegraphics[width=\columnwidth]{flowchart_model.pdf}
    \caption{经典 Mesh 接触模型、传统一阶 SDF 接触模型与本文提出的具备 Hessian 曲率感知的二阶 SDF 接触模型的管线架构对比。}
    \label{fig:flowchart}
\end{figure}

\subsection{基础 NSC 局部一阶预测缺陷模型}

在基于一阶线性预测的 LCP 接触处理中，为了在整个时间步 $\Delta t$ 后阻止两物体表面过度穿透，下一时刻速度 $\mathbf{v}^{l+1}$ 需满足局部 Signorini 非穿透条件的一阶线性化：

\begin{equation}
    \Phi(\mathbf{q}^{l+1}) \approx \Phi(\mathbf{q}^l) + \Delta t \, \mathbf{n}^T \mathbf{v}^{l+1} \geq 0
\end{equation}

其中，$\Phi(\mathbf{q}^l)$ 为当前步的接触间隙，$\mathbf{n}$ 为接触法向。对于基于 SDF 的点接触模型，该法向可由距离场梯度给出。为了与后续二阶曲率构造保持一致，本文统一采用单位法向
\begin{equation}
    \mathbf{n} = \frac{\nabla \Phi}{\|\nabla \Phi\|}.
\end{equation}

在传统一阶模型中，当时间步较大或接触点沿高曲率边界发生快速切向滑移时，局部切平面外推会偏离真实表面流形，从而使 RHS 中的间隙预测出现系统性失真。这类误差会被非光滑互补主方程放大为非物理冲量，最终表现为速度曲线中的高频毛刺与脉冲。

\subsection{接触流形 Weingarten 映射与曲率解析}

为缓解一阶切平面模型的几何失配，本文将接触间隙的局部展开提升到二阶。实现上，我们并不直接对标量距离场 $\Phi$ 求 Hessian，而是先构造单位法向场
\begin{equation}
    \mathbf{n}(\mathbf{x}) = \frac{\nabla \Phi(\mathbf{x})}{\|\nabla \Phi(\mathbf{x})\|},
\end{equation}
再对其做中心差分，以获得局部曲率张量
\begin{equation}
    \mathbf{W}(\mathbf{x}) \approx \nabla \mathbf{n}(\mathbf{x}).
\end{equation}

因此，本文所使用的二阶张量本质上是离散 SDF 法向流形上的 Weingarten 映射近似，而不是直接对 $\Phi$ 标量场求原始二阶导。对于当前时间步的局部相对位移 $\mathbf{s} = \Delta t \, \mathbf{v}_{rel}$，接触间隙在二阶上可写为
\begin{equation}
    \Phi(\mathbf{q}^{l+1}) \approx \Phi(\mathbf{q}^{l}) + \mathbf{s}^T \mathbf{n} + \frac{1}{2}\mathbf{s}^T \mathbf{W}\mathbf{s}.
\end{equation}

进一步采用当前步相对速度的半隐式近似，可得到二阶曲率补偿量
\begin{equation}
    C_{curve} = \frac{1}{2} (\Delta t)^2 \, \mathbf{v}_{rel}^T \mathbf{W} \mathbf{v}_{rel}.
\end{equation}

\subsection{LCP 主方程右端项的曲率前馈注入}

直接改写互补问题左端的动力学矩阵会破坏底层 NSC/LCP 结构的稳定性与可复用性。为保持互补主方程的形式不变，本文采用寄生式 RHS 注入策略，将二阶几何信息以前馈偏置的形式写入有效接触间隙：
\begin{equation}
    \Phi_{eff} = \Phi + C_{curve}.
\end{equation}

由此，原始一阶非穿透条件可重写为
\begin{equation}
    \Delta t \, \mathbf{n}^T \mathbf{v}^{l+1} \geq - \Phi_{eff}.
\end{equation}

也即，在不改写底层互补模块整体结构的前提下，通过修正 RHS 中的接触间隙，向 LCP/NSC 主方程注入局部二阶几何预测信息。该处理保留了传统非光滑接触框架的主干求解流程，同时将高曲率表面上的运动学提前量并入接触判定。

\subsection{接触类型感知的局部接触表示}

本文并不假设单一接触代理能够同样适用于所有接触场景，而是根据局部接触的几何与运动学特征采用两类不同的接触表示。对于工业凸轮、偏心圆盘-滚子以及 Onset Stress 等光滑连续滑移接触，本文采用 \texttt{SlidingPatch} 表示：接触集合在时间上保持 persistent manifold 身份，并通过 representative query、局部拟合与候选筛选来维持连续的 patch 级几何更新。对于 Twin Gear 一类紧凑主导点及非光滑奇异边界接触，本文采用 \texttt{CompactDiscrete} 表示：patch 的目标主要是进行去重、法向稳健化与主导几何保留，而不再追求连续滑移流形的时间追踪。

这一分流的核心思想在于：高阶曲率项本身并不保证对所有接触都产生正收益，它只有在与接触类型相匹配的局部流形表示中才能发挥稳定作用。对光滑连续接触而言，局部流形的持续追踪是二阶前馈获得净增益的前提；而对非光滑或分支切换显著的接触而言，更保守的紧凑离散表示反而有助于抑制由高阶几何噪声引起的非物理脉冲。

图 \ref{fig:contact_regime_flow} 进一步总结了本文在运行时的三阶段处理路径：系统首先执行候选样本预处理，并可在样本侧启用 OBB/BVH 分层 broad-phase 以降低候选查询开销；随后在局部流形构造阶段根据接触的滑移连续性与主导几何特征选择 \texttt{SlidingPatch} 或 \texttt{CompactDiscrete}；最后将修正后的有效间隙统一写回 NSC/LCP 主方程。

\begin{figure}[htbp]
    \centering
    \includegraphics[width=\columnwidth]{flowchart_contact_regime.pdf}
    \caption{Contact-type-aware routing used in the proposed framework. The runtime pipeline is organized into candidate preprocessing, regime-aware manifold construction, and constraint injection. Candidate samples may first pass through an optional sample-side OBB/BVH broad-phase; smooth continuous contacts are then processed by \texttt{SlidingPatch}, while compact or singular contacts are processed by \texttt{CompactDiscrete}, and both branches finally inject effective gaps into the same NSC/LCP solve chain.}
    \label{fig:contact_regime_flow}
\end{figure}

\subsection{Patch 级接触聚类与局部流形约化}

由于离散采样会在同一局部接触区产生多个近邻样本，若逐点直接将所有 SDF 查询结果送入互补模块，法向噪声与重复接触点会显著放大数值抖动。为此，本文在活动接触集合上引入 patch 级聚类策略。

首先，对所有候选活动样本按穿透深度排序；随后依据空间邻近性与法向一致性进行聚类。对任意候选样本 $\mathbf{x}_i$ 与已建立 patch 的代表样本 $\mathbf{x}_k$，若满足
\begin{equation}
    \|\mathbf{x}_i - \mathbf{x}_k\| < r_{cluster},
\end{equation}
\begin{equation}
    \mathbf{n}_i \cdot \mathbf{n}_k > \cos \theta_{cluster},
\end{equation}
则将其并入同一 patch。其中，$r_{cluster}$ 为空间半径阈值，$\theta_{cluster}$ 为法向夹角阈值。

在每个 patch 内，法向向量采用样本平均并重新归一化：
\begin{equation}
    \bar{\mathbf{n}}_k =
    \frac{\sum_{i \in M_k} \mathbf{n}_i}
         {\left\|\sum_{i \in M_k} \mathbf{n}_i\right\|}.
\end{equation}
但 patch 的主几何量并不做完全平均。对于穿透深度 $\Phi$、梯度以及曲率张量 $\mathbf{W}$，本文保留该 patch 中 deepest sample 的取值，即保留最深穿透样本所对应的局部几何状态。于是每个 patch 最终由
\begin{equation}
    \left(\bar{\mathbf{n}}_k, \Phi_k^{deep}, \mathbf{W}_k^{deep}\right)
\end{equation}
共同表征，并进一步参与互补模块中的接触约束组装。

这种处理兼顾了两个目标：一方面通过法向平均降低离散点接触带来的方向抖动；另一方面通过保留 deepest sample 的 $\Phi/Hessian$，避免对主接触几何进行过度平滑，从而保留局部曲率补偿对真实接触主导点的响应能力。
